%\clearpage

\section{Summary of changes included in version 8 of T2K-TN-141 (released on July 09, 2013)}

This section is designed to allow readers to keep track of the changes between versions of T2K-TN-141.
Table, figure, equation, footnote or citation numbers shown in black refer to version of 7,
{\color{red}whereas the numbers shown in red refer to version of 8}.
Both numbers are shown where possible.

\begin{itemize}

\item
  Information on, and studies involving the assumption of the second octant
  ($\sin^2(\theta_{23}) \geq 0.5$) have been added. All additional fits have been done with the
  '2+N' fitter, and use the Run 1+2+3 dataset only.
  \begin{itemize}
  \item Discussion has been added to Section 2.6 {\color{red}(2.6)} on why the
    choice of octant assumption makes a difference to the prediction at SuperK.
  \item The effect of the octant choice on the SuperK reconstructed energy spectrum in shown in Fig.\ {\color{red}12}.
  \item The effect of the octant choice on the expected number of SuperK events in shown in Fig.\ {\color{red}13}.
  \item The addition of the result '2+N' fit result to the Run 1+2+3 dataset assuming the second octant is shown in Section 5.1 {\color{red}(5.1)},
    and in Table 13 {\color{red}(13)}.
  \item The pulls of the '2+N' fit result assuming the second octant is shown in Fig.\ {\color{red}34}.
  \item The goodness of fit of the '2+N' fit result assuming the second octant is shown in Fig.\ {\color{red}41}.
  \item The best-fit spectrum of the '2+N' fit result assuming the second octant is shown in Fig.\ {\color{red}44} for the fine binning scheme,
    Fig.\ {\color{red}52} for the coarse binning scheme, and Fig.\ {\color{red}53} for a breakdown by MC component prediction.
  \item A comparison of the best-fit spectra of the '2+N' fits assuming the first and second octants is shown in Figs.\ {\color{red}56} and {\color{red}57},
    for the fine and coarse binning schemes respectively.
  \item The 1D constant $\Delta \chi^2$ confidence limits for '2+N' fits, assuming the second octant have been added to Section 5.4 {\color{red}(5.4)}.
    The corresponding 1D $\Delta \chi^2$ plots have been added to Figs.\ 51-52 {\color{red}(60-61)}.
  \item The $\Delta \chi^2$ surface for the '2+N' fit assuming the second octant is shown in Fig.\ {\color{red}63}.
  \item The fractional difference between $\Delta \chi^2$ surfaces for both octant assumptions, are shown in Fig.\ {\color{red}73}.
  \item The contour for the '2+N' fit assuming the second octant, made using the constant $\Delta \chi^2$ method is shown in Fig.\ {\color{red}79}.
  \item The contour for the '2+N' fit assuming the second octant, made using the FC method is shown in Fig.\ {\color{red}86}.
  \item Section {\color{red}5.6.1} has been added, comparing '2+N' fit contours, for both octant assumptions.
  \item A comparison between octant assumptions for '2+N' fits to the Run 1+2+3 dataset is shown in Figs.\ {\color{red}88} and {\color{red}89}
    for the constant $\Delta \chi^2$ and FC methods respectively.
  \item A comparison of contours between both octant assumptions, converted into $\sin^2\theta_{23}$-$|\Delta m^2_{32}|$ is shown in Fig.\ {\color{red}90}.
    Also shown is the result of a fit to sin$^2\theta_{23}$ directly. The best-fit points in this space are listed in Tab.\ {\color{red}14}.
  \item A comparison between the constant $\Delta \chi^2$ and FC methods for the '2+N' assuming the second octant is shown in Fig.\ {\color{red}95}.
  \item A comparison between the '2+N' contours assuming the second octant with external data is shown in Fig.\ {\color{red}98}.
  \item A comparison between the '2+N' FC contours, for both octant assumptions, with external data is shown in Fig.\ {\color{red}100}.
  \item Added the result of the '2+N' best-fit point in the second octant and the corresponding 1D limits to both the Abstract, and Sec.\ 6 {\color{red}(6)}.
  \end{itemize}
  Note that the FC contours will be updated as more toy MC statistics become available for making the $\Delta\chi^2_{\textrm{crit}}$ surface.
  
\end{itemize}
